% Generated draft: requires mathematical and source review.
\documentclass[11pt]{article}
\usepackage[margin=25mm]{geometry}
\usepackage{fontspec}
\setmainfont{Latin Modern Roman}
\usepackage{amsmath,amssymb,mathtools,unicode-math}
\setmathfont{Latin Modern Math}
\usepackage{enumitem,needspace,xurl,hyperref}
\setlength{\parindent}{0pt}
\setlength{\parskip}{5pt}
\setlength{\emergencystretch}{3em}
\newcommand{\N}{\mathbb{N}}
\newcommand{\Z}{\mathbb{Z}}
\newcommand{\Q}{\mathbb{Q}}
\newcommand{\R}{\mathbb{R}}
\newcommand{\C}{\mathbb{C}}
\newcommand{\F}{\mathbb{F}}
\newcommand{\A}{\mathbb{A}}
\newcommand{\PP}{\mathbb{P}}
\newcommand{\E}{\mathbb{E}}
\newcommand{\eps}{\varepsilon}
\newcommand{\dd}{\,\mathrm{d}}
\begin{document}
\section*{UnsolvedMath ID: 30002065}
\textbf{Source problem code:} catalogue entry
\subsection{Definitions and mathematical statement}
Let \(\Delta^d\) denote a \(d\)-dimensional simplex and let \(\partial\Delta^d\) be its boundary simplicial complex. For each nonnegative integer \(n\), define
\[
K_{d,n}=\operatorname{sd}^{\,n}(\partial\Delta^d),
\qquad K_{d,0}=\partial\Delta^d,
\]
where \(\operatorname{sd}\) denotes barycentric subdivision. The source does not specify the intended range of dimensions \(d\), including the treatment of low-dimensional or empty boundaries.

\textbf{Unresolved formulation data.} The record does not define which Morse functions are allowed, how a random function is sampled, or what numerical invariant is meant by its ``complexity (complicatedness).'' The supplied literature references suggest a discrete Morse setting, but they do not, in the supplied text, specify a sampling algorithm or complexity functional. Thus the following is a conditional formulation, rather than a uniquely specified probabilistic problem.

For each admissible pair \((d,n)\), suppose that a measurable space \(\mathcal M_{d,n}\) of the intended Morse functions on \(K_{d,n}\), a probability measure \(\mu_{d,n}\), and a measurable complexity functional \(C_{d,n}\) have been specified. Whenever the expectation is defined, write
\[
a_{d,n}=\mathbb E_{f\sim\mu_{d,n}}\!\left[C_{d,n}(f)\right].
\]
\textbf{Source question.} Determine the expected complexity \(a_{d,0}\) of a random Morse function on the boundary of a \(d\)-simplex, and determine how \(a_{d,n}\) behaves under repeated barycentric subdivision, including its growth as \(n\) increases. This question requires the intended class of functions, probability laws, complexity functionals, and dimension range to be recovered before a precise unconditional statement can be given.

\textbf{Supplied status qualification.} The registry labels the problem partially solved. Its literature assessment dated 2026-08-21 states that the repeated-subdivision subquestion has an exponential lower-bound answer, while the expected complexity on simplex boundaries is not determined in general. The supplied record does not state the hypotheses, dimensions, random model, complexity invariant, constants, or quantifiers for that lower bound. Accordingly, that assessment is retained as supplied dated curation and is not asserted here as a verified theorem about the conditional quantities \(a_{d,n}\).
\subsection{Short English statement}
The question asks for the expected complexity of a random Morse function on a simplex boundary and for its behavior after repeated barycentric subdivision. A precise expectation requires a definition of the Morse functions, their probability distribution, and their complexity. These are missing from the supplied record. The supplied dated assessment labels the problem partially solved and reports an exponential lower bound for the subdivision subquestion; its precise scope requires review.
\subsection{Sources}
\begin{itemize}
\item Supplied local UnsolvedMath registry record, permanent ID \(30002065\), title ``Complexity of Random Morse Functions on Simplex Boundaries,'' in the Oberwolfach Reports Open Problems collection. Exact source problem code: \texttt{catalogue entry}. The registry gives proposed year 2012 and names no proposer.
\item Supplied source citation: \emph{Triangulations (2013)}, Oberwolfach Reports, DOI 10.4171/owr/2012/24, pp.~1405--1486. \url{https://doi.org/10.4171/owr/2012/24}. The background describes these reports as documenting mathematical workshops and their open-problem discussions. The supplied year label and the year appearing in the DOI require bibliographic reconciliation; no problem-specific page is supplied.
\item Supplied literature source: Adiprasito--Benedetti--Lutz. \url{https://doi.org/10.1007/s00454-017-9860-4}.
\item Supplied literature source: ``random discrete Morse theory.'' \url{https://arxiv.org/abs/1303.6422}.
\item Supplied literature source: ``original OWR report.'' \url{https://doi.org/10.4171/owr/2012/24}.
\item Literature assessment date: \textbf{2026-08-21}. This is supplied dated curation derived from the supplied OWR curation, not independently verified. It reports an exponential lower-bound answer for the repeated-barycentric-subdivision subquestion and says that the expected complexity for random Morse functions on simplex boundaries is not determined in general. The source explicitly calls for independent verification of time-sensitive status and cited claims.
\item The registry supplies statement status \texttt{corrected\_verified} and describes a manual semantic, typographic, and LaTeX review of a meaning-preserving cleanup. This is supplied editorial provenance, not independent verification of the mathematical formulation or literature assessment.
\end{itemize}
\end{document}
